%=====================================================================
\chapter{Extending the Platform: Your First Module}\label{ch:modules}
%=====================================================================
\locator{5}{Part V --- Running and Extending the Shop}

\begin{outcomes}
By the end of this chapter you will be able to:
\begin{tight}
  \item \textbf{Explain} what an extension is made of, and \textbf{state} why
        each part exists.
  \item \textbf{Write} a storefront module and a shipping method that add
        behaviour without editing anything the platform ships.
  \item \textbf{Install} an extension the way the marketplace installer does,
        in three steps you can read.
  \item \textbf{Measure} what a module costs a page, and \textbf{account} for
        every query it adds.
  \item \textbf{Judge} when to extend, when to configure, and when to refuse.
\end{tight}
\end{outcomes}

\begin{prereq}
\chref{ux}: themes, templates and overrides. \chref{shipping}: the shipping
quote interface and tax classes. \chref{catalogue}: counting the queries a page
runs. Appendix~B for enough PHP and Twig to read the files.
\end{prereq}

\begin{keyterms}
extension $\bullet$ module $\bullet$ extension point $\bullet$ namespace
$\bullet$ autoloading $\bullet$ manifest $\bullet$ controller $\bullet$ model
$\bullet$ language file $\bullet$ module instance $\bullet$ layout position
$\bullet$ upgrade safety $\bullet$ permission check
\end{keyterms}

\section{Why Not Just Edit the File}

Every feature in this chapter could be added faster by opening the shop's own
code and changing it. The reason not to is one sentence long, and it arrives
about eighteen months later: \textbf{the next upgrade overwrites your file, and
nothing tells you which ones you changed.}

\begin{conceptbox}[The three ways to change what a shop does, in order of
preference]
\begin{tightnum}
  \item \textbf{Configure it.} If a setting exists, use it.
        \chref{shipping}'s tax rules and \chref{dashboard}'s widgets are
        settings, and a setting survives everything.
  \item \textbf{Extend it.} Add a directory under \texttt{extension/} that the
        platform loads. Nothing it ships is touched, the upgrade leaves your
        directory alone, and \texttt{git status} tells you exactly what you
        added.
  \item \textbf{Edit it.} Only when the first two cannot reach, and then with
        the change recorded somewhere an upgrade will make you re-read.
\end{tightnum}

\smallskip
\textbf{Most commissioned work on this platform is the second.} That is what
the department's title adds to HEC's bare ``E-Commerce'', and it is the
difference between configuring a shop and developing one.
\end{conceptbox}

\section{What an Extension Is Made Of}

\begin{definitionbox}[Extension, Module, Instance]
An \term{extension} is a directory under \texttt{extension/} whose internal
shape mirrors the shop's own --- an \texttt{admin/} half, a \texttt{catalog/}
half, and a manifest. It may provide several things at once: this chapter's
provides a module and a shipping method.

\smallskip
A \term{module} is a block of storefront output that a layout can place. It
has two controllers: one that renders it for a customer, one that lets the
shopkeeper configure it.

\smallskip
An \term{instance} is one placement of a module, with its own settings. The
same module may appear three times on a page with three different
configurations, which is why its render method is handed its settings rather
than reading them from the global configuration.
\end{definitionbox}

\dsfig{%
\begin{tikzpicture}[
  dir/.style={rounded corners=2pt, draw=#1, line width=0.9pt, fill=#1!8,
              text=#1!60!black, font=\bfseries\scriptsize, align=left,
              text width=7.0cm, minimum height=0.62cm, inner sep=3pt},
  note/.style={font=\scriptsize, text=neutral, align=left, text width=5.2cm,
               anchor=west}
]
  \node[dir=purplehl]  (m) at (0,3.0)  {extension/ecadmod/install.json};
  \node[dir=secondary] (a) at (0,2.1)  {\ \ admin/controller/module/lowstock.php};
  \node[dir=secondary] (t) at (0,1.2)  {\ \ admin/view/template/\dots};
  \node[dir=greenhl]   (c) at (0,0.3)  {\ \ catalog/controller/module/lowstock.php};
  \node[dir=greenhl]   (v) at (0,-0.6) {\ \ catalog/view/template/\dots};
  \node[dir=neutral]   (l) at (0,-1.5) {\ \ catalog/language/en-gb/\dots};

  \node[note, text=purplehl] at (7.3,3.0)
       {Who wrote it, and what it is. Without this row the admin panel cannot
        see the extension, however complete the files are.};
  \node[note, text=secondary] at (7.3,1.65)
       {The shopkeeper's half: a form, validation, and a save action that
        checks a permission before it does anything.};
  \node[note, text=greenhl] at (7.3,-0.15)
       {The customer's half: one method, handed its settings, returning a
        string.};
  \node[note, text=neutral] at (7.3,-1.5)
       {Every visible word. A shop translated later needs no code change.};
\end{tikzpicture}}%
{The twelve files of this chapter's extension. The namespace of every class is
its path, which is how the loader finds it and the first thing to get wrong.}
{extensionshape}

\begin{alertbox}[The namespace is the path, and the error message is unhelpful]
The class

\smallskip
\hspace*{1em}\texttt{Opencart\textbackslash Catalog\textbackslash
Controller\textbackslash Extension\textbackslash Ecadmod\textbackslash
Module}

\smallskip
\noindent lives in
\texttt{extension/ecadmod/catalog/controller/module/}. Get a single segment
wrong --- the wrong case, a missing \texttt{Extension}, the folder named
differently from the namespace --- and the loader reports that the class does
not exist.

\smallskip
Which is true, and tells you nothing about why. \textbf{When an extension does
not load, check the namespace against the path before anything else}; it is
the cause more often than every other possibility combined.
\end{alertbox}

\section{The Storefront Half, in Full}

The whole customer-facing side of the module is one method. It is worth reading
in full, because almost nothing in it is specific to this module.

\begin{lstlisting}[language=php]
namespace Opencart\Catalog\Controller\Extension\Ecadmod\Module;

class LowStock extends \Opencart\System\Engine\Controller {
    public function index(array $setting): string {
        $this->load->language('extension/ecadmod/module/lowstock');

        $data['heading_title'] = $this->language->get('heading_title');
        $threshold = (int)($setting['threshold'] ?? 5);
        $limit     = (int)($setting['limit'] ?? 4);

        $query = $this->db->query("SELECT p.product_id,
            p.quantity, pd.name FROM ...
            WHERE p.status = '1' AND p.quantity > '0'
              AND p.quantity <= '" . (int)$threshold . "'
            ORDER BY p.quantity ASC LIMIT " . (int)$limit);
\end{lstlisting}

\noindent and then, in the same method, the rows become view data and the
view becomes a string:

\begin{lstlisting}[language=php]
        $data['products'] = [];

        foreach ($query->rows as $row) {
            $data['products'][] = [
                'name'     => $row['name'],
                'quantity' => $row['quantity'],
                'href'     => $this->url->link('product/product',
                    'language=' . $this->config->get('config_language')
                    . '&product_id=' . $row['product_id'])
            ];
        }

        if (!$data['products']) { return ''; }

        return $this->load->view(
            'extension/ecadmod/module/lowstock', $data);
    }
}
\end{lstlisting}

\begin{conceptbox}[Four decisions in thirty lines]
\begin{tightnum}
  \item \textbf{It takes \texttt{\$setting} as an argument} rather than reading
        \texttt{\$this->config}. That is what allows the module to be placed
        twice with different thresholds.
  \item \textbf{It returns a string and never prints.} The layout decides where
        the string goes; a module that echoes has taken a decision that is not
        its own.
  \item \textbf{It returns \texttt{''} when there is nothing to say.} An empty
        panel with a heading is worse than no panel.
  \item \textbf{Every integer is cast.} \texttt{(int)\$threshold} inside a
        query string is the difference between a parameter and an injection
        --- and \chref{security} showed that the platform is careful about this
        everywhere, which means your extension is now the weakest part.
\end{tightnum}
\end{conceptbox}

\section{Installing It, in Three Steps You Can Read}

The marketplace installer is a button. What it does is three things, and doing
them by hand once is the fastest way to understand why an extension that looks
installed sometimes is not.

\begin{examplebox}[What \texttt{code/m19\_modules.py --install} performs]
\begin{tightnum}
  \item \textbf{Copy the files} to \texttt{extension/ecadmod/} and make them
        readable by the web server. Twelve files.
  \item \textbf{Insert the manifest row} into \texttt{oc\_extension\_install}
        --- name, author, version, link. \textbf{Without this the admin panel
        cannot see the extension at all}, however complete the files are. This
        is the step people miss, and the symptom is an extension that is
        plainly on disk and absent from every list.
  \item \textbf{Insert one row per thing provided} into
        \texttt{oc\_extension}: \texttt{(ecadmod, module, lowstock)} and
        \texttt{(ecadmod, shipping, zonerate)}. This is what makes each appear
        in its own list rather than all of them in one.
\end{tightnum}

\smallskip
Uninstalling is the same three in reverse, and the script does that too ---
because an extension you cannot remove cleanly is one you cannot test.
\end{examplebox}

\section{What a Module Costs}

The objection to building features as modules is that they must be slow. It is
worth measuring rather than assuming, in both directions.

\begin{measured}[The same home page with 0, 1, 3 and 5 instances]
From \texttt{code/m19\_modules.py --measure} against the stack in Appendix~A.
Each instance is a separate placement of the same module on the home layout.

\rowcolors{2}{white}{lightbg}
\begin{longtable}{@{}L{2.4cm}L{2.4cm}L{2.4cm}L{3.0cm}L{3.2cm}@{}}
\toprule
\rowcolor{primary!15}
\textbf{Instances} & \textbf{Page} & \textbf{Queries} & \textbf{ms each} &
\textbf{queries each} \\
\midrule
\endhead
0 & 76\,ms & 117 & --- & --- \\
1 & 82\,ms & 121 & 5.8 & \textbf{4.0} \\
3 & 90\,ms & 129 & 4.7 & \textbf{4.0} \\
5 & 96\,ms & 137 & 4.0 & \textbf{4.0} \\
\bottomrule
\end{longtable}

\smallskip
\textbf{The cost is linear and small}: about 4\,ms and exactly four queries per
instance. Five instances added 20\,ms to a 76\,ms page. A module is not
expensive.

\smallskip
\textbf{But the module contains \emph{one} query.} The other three are the
framework's, and reading them from the slow log names each one:

\smallskip
\begin{tabular}{@{}L{7.4cm}L{6.2cm}@{}}
\texttt{SELECT * FROM oc\_module WHERE module\_id=16} & this instance's
settings \\
\texttt{SELECT * FROM oc\_translation WHERE route=\dots} & its translations \\
\texttt{SELECT p.product\_id, p.quantity \dots} & \textbf{the one I wrote} \\
\texttt{SELECT * FROM oc\_theme WHERE route=\dots} & its template override \\
\end{tabular}

\smallskip
\textbf{Three quarters of what a module costs is the extension mechanism
itself}, spent finding your settings, your words and your template. That is the
price of the arrangement that lets you place the module twice, translate it
without touching code, and override its template from a theme.

\smallskip
It is also the explanation for \chref{catalogue}'s 144 queries on a category
page, which was never explained at the time: a page is mostly modules, and
every one of them pays this toll.
\end{measured}

\begin{conceptbox}[When to extend, and when to refuse]
\textbf{Extend when} the behaviour is yours, the platform offers a point for
it, and you can state what the shop would do without it.

\smallskip
\textbf{Configure instead when} a setting already reaches. Rewriting a setting
as a module is the commonest waste in commissioned work.

\smallskip
\textbf{Refuse when} the request is for something the platform deliberately does
not do, and the reason it does not is a good one. ``Store the card number so
customers do not have to retype it'' is a module somebody will ask you to
write, and \chref{gateway} is the reason to say no --- in writing, with the
consequence named.
\end{conceptbox}

\begin{pitfall}
\begin{tight}
  \item \textbf{Editing a shipped file} because it is faster today.
  \item \textbf{A namespace that does not match the path.} The error says the
        class does not exist, which is unhelpfully true.
  \item \textbf{Forgetting the \texttt{oc\_extension\_install} row}, which
        leaves an extension on disk and invisible.
  \item \textbf{Reading \texttt{\$this->config} in a module's render method}
        instead of its \texttt{\$setting} argument, so it cannot be placed
        twice.
  \item \textbf{Omitting the permission check in a save action.} The form is
        only shown to an administrator; the action is a URL.
  \item \textbf{Interpolating an uncast value into SQL.} The platform is
        careful everywhere else, so your extension becomes the weakest part.
  \item \textbf{Returning a tax-inclusive cost from a shipping quote}, which
        the totals pipeline then taxes again (\chref{shipping}).
  \item \textbf{Hard-coding visible text} in a controller or template rather
        than a language file.
\end{tight}
\end{pitfall}

\begin{lab}[Write, install, measure, and then break it]
\begin{lstlisting}[language=bash]
# --- 1. Read the twelve files before running anything ---------------
cd code
find modules/ecadmod -type f | sort
# Match each class's namespace to its path. Say the rule out loud.

# --- 2. Install it the way the marketplace would --------------------
python m19_modules.py --install

# --- 3. Place it, and look --------------------------------------
python m19_modules.py --place 1
# Open the storefront. Nothing shows -- the demo products all have
# high stock. Fix that by raising the threshold, and say why an
# empty module returning '' is better than an empty panel.

# --- 4. Predict the cost, then measure it ---------------------------
# How many queries does ONE instance add? Commit to a number.
python m19_modules.py --measure

# --- 5. Account for all four ----------------------------------------
docker compose exec db mariadb -uroot -prootpw -e \
  "SELECT LEFT(CONVERT(sql_text USING utf8),80) FROM mysql.slow_log;"
# One is yours. Name the other three and say what each is for.

# --- 6. Break the namespace on purpose ------------------------------
# Rename the class to Lowstock (small s) and reload. Read the error.
# This is the error you will see most often; meet it deliberately now.

# --- 7. Configure your shipping method -------------------------------
# Admin: Extensions > Extensions > Shipping > Zone Rate > Edit.
# Rates one per line as zone_id:cost. Place an order and check the
# quote, then check oc_order_total for how the tax was applied.

# --- 8. Add the free-delivery threshold and test both sides ---------
# One order below it, one above. Verify from oc_order_total, not
# from the screen.

# --- 9. Remove it completely ----------------------------------------
python m19_modules.py --remove
# Then confirm the shop is exactly as it was: no files, no rows,
# no settings. An extension you cannot remove is one you cannot test.
\end{lstlisting}

\smallskip
\textbf{What to notice.} Step 5 is the chapter. Your module is one query; it
costs four. Knowing which three are the mechanism's is what lets you judge
whether a page with twenty modules is badly written or merely busy.
\end{lab}

\begin{checkpoint}
\begin{tightnum}
  \item Five instances of a module added 20\,ms and 20 queries to a 76\,ms
        page. A colleague concludes modules are expensive and proposes editing
        the theme instead. Give your answer, with the per-instance breakdown,
        and say what you would measure before agreeing to anything.
  \item A module's render method reads its threshold from
        \texttt{\$this->config->get('module\_lowstock\_threshold')} instead of
        from its \texttt{\$setting} argument. The shop works. Describe exactly
        when it stops working and what the shopkeeper will report.
  \item You are asked to add a field to the checkout that the platform does not
        offer. Set out how you would decide between configuring, extending and
        editing, what you would need to know first, and what you would write
        down whichever you chose.
\end{tightnum}
\end{checkpoint}

\begin{chaptersummary}
\begin{tight}
  \item Configure first, extend second, edit last. An extension is a directory
        the upgrade leaves alone; an edit is a file the upgrade overwrites
        without telling you.
  \item An extension mirrors the shop's shape: an admin half, a catalog half, a
        manifest, and language files for every visible word.
  \item The namespace of every class is its path. A mismatch reports that the
        class does not exist, and is the commonest cause by a wide margin.
  \item A module's render method takes its settings as an argument, returns a
        string, returns \texttt{''} when it has nothing to say, and casts every
        value it puts in a query.
  \item Installing is three steps: the files, a manifest row in
        \texttt{oc\_extension\_install}, and one row per provided thing in
        \texttt{oc\_extension}. Miss the second and the extension is on disk
        and invisible.
  \item Measured: 4\,ms and exactly 4 queries per module instance; five
        instances added 20\,ms to a 76\,ms page. Modules are cheap.
  \item Measured: of those four queries, \emph{one} is the module's. The other
        three find its settings, its translations and its template --- the
        price of being placeable, translatable and overridable, and the reason
        a category page runs 144 queries.
  \item A shipping method's whole interface is one method that is handed an
        address and returns a quote without tax, or nothing at all.
  \item Refusing is a professional act. ``Store the card number'' is a module
        somebody will ask for; decline it in writing, with the consequence
        named.
\end{tight}
\end{chaptersummary}

\begin{reviewq}
  \item \textbf{(MCQ)} One instance of the measured module added how many
        database queries? (a)~1; (b)~2; (c)~4; (d)~it varies with the
        catalogue.
  \item \textbf{(MCQ)} An extension whose files are all present but which
        appears in no admin list is most likely missing: (a)~a language file;
        (b)~its \texttt{oc\_extension\_install} row; (c)~file permissions;
        (d)~a template.
  \item \textbf{(Short)} Explain why a module's render method is handed its
        settings instead of reading the global configuration.
  \item \textbf{(Short)} Give the rule relating a class's namespace to its file
        path, and the symptom of breaking it.
  \item \textbf{(Short)} Why must a shipping quote return a cost without tax?
        Refer to \chref{shipping}'s totals pipeline.
  \item \textbf{(Applied)} You are asked for a module showing ``customers who
        bought this also bought''. Specify it: what it queries, where its
        settings live, how many queries it will add per page, how you would
        keep that number from growing with the catalogue, and what you would
        measure before and after.
  \item \textbf{(Applied)} A client wants a field on the product page that
        their staff fill in and customers see. Compare doing this with a
        product attribute (configuration), a module (extension) and a template
        edit. Give the work, the upgrade risk and the maintenance burden of
        each, and recommend one with your reasons.
\end{reviewq}
